\section{Conclusion and Future Work}
\label{sec:conclusion}

We have argued that deciding which open educational resources to adopt needs a shared instrument that is anchored, two-part, reproducible and validated, and that existing rubrics, as surveyed in the literature, meet these requirements only partly. We proposed one candidate, the OpenBook QR Score, a ten-metric, 100-point instrument whose report keeps quality and reusability visible as QR 76 (Q 37 / R 39) and whose decision rule guards against a high total hiding a failed metric. An illustration on ten widely used textbooks, scored in two passes of a few minutes of wall-clock time each, showed the instrument's value (it points to concrete fixes, and the two passes agreed within one point on 88 percent of metric scores) and its weaknesses (descriptors for restrictive licenses are loose, practice with feedback is unstable across chapters, and all scores come from one model family).

Four pieces of future work follow directly. First, tighten descriptors where the illustration exposed drift, starting with licenses and learning outcomes. Second, run the human and agent agreement studies of Section~\ref{sec:validation}, including at least two independent human reviewers per resource and a sample of several chapters per book. Third, test whether scores predict learner outcomes in controlled comparisons of an original resource and an improved version. Fourth, publish scores with their evidence and version in an open catalog, so that the community can dispute and improve them. Until those studies are done, every OpenBook QR Score should be labeled provisional.
